\section{Discussion of Findings: Implications For Smart Building Research}
\label{sec:ch4-discussions}
Our findings include seven \emph{topics (T1-T7)}, eight \emph{themes (T1-T8)}, and five \emph{design opportunities (DO1-DO5)}. While \emph{topics} served as an organising scaffold, here we focus on how the \emph{themes} and \emph{opportunities} extend HBI research: \emph{thematically}, by expanding the experiential and ecological vocabulary of interaction, and \emph{generatively}, by proposing new design directions and technological opportunities. We then share messages for a multisensory–multidimensional (MS–MD) approach to biophilic interaction. Finally, we discuss tensions in applying biophilic design and study limitations.

\subsection{Thematic Contributions: Expanding the Vocabulary of HBI} 
Our findings contribute to HBI by expanding its design vocabulary---re-framing smart buildings as \emph{ecological interfaces} that surface natural rhythms and multispecies presence \textbf{(T1, T3, T5)}. While prior HBI work has largely prioritised comfort and behavioural control~\cite{alavi2017comfort, nguyen2024adaptive}, ecological awareness and more-than-human perspectives remain under-explored. Through our themes, we surface terms such as \emph{multispecies flourishing}, \emph{biophilic (dis)comfort}, and \emph{ecological interface} to capture this shift. Second, we foreground \emph{multisensory, somatic engagement} in interaction design. Experts envisioned systems that respond to the body through movement, rhythm, and material sensitivity---beyond visual dashboards or app-based feedback \textbf{(T2, T4, T7, T8)}. While soma design has highlighted embodied experience in HCI~\cite{hook2018designing}, our findings extend these insights into architectural space, suggesting how buildings themselves may become expressive, responsive participants in biophilic interaction. Finally, our themes emphasise the \emph{social and cultural responsibility} of smart environments. Although HBI has begun to address equity and privacy~\cite{mitchell2020human}, it often overlooks how cultural values, neurodiversity, and uneven access to nature shape lived experience. Biophilic design offers a way to support diverse bodies, histories, and sensory needs—enabling justice-oriented smart environments~\cite{chordia2024social}.

\subsection{Generative Opportunities: Biophilic Interaction Through Technology}
Our design opportunities contribute to HBI by reimagining how technology can support biophilic interaction. These contributions open new conceptual and material directions for IxD in the built environment. First, we reposition \emph{sensing} from a behavioural optimisation~\cite{chhaglani2022flowsense} mechanism to a medium for ecological expression. Speculative deployments of stochastic IoT and generative AI render light, sound, and air movement as expressive materials---carrying the rhythms of weather, growth, or animal migration \textbf{(DO1, DO5)}. This extends sensing beyond data extraction into ambient ecological storytelling. Second, shifting from traditional HCI perspectives, we build on soma design to propose \emph{embodied interactions at architectural scale}. Haptic panels, responsive facades, and AR overlays attuned to movement and physiology foster bodily and interpretive engagement with biophilic environments \textbf{(DO2)}. These approaches challenge the dominance of reactive, screen-based interfaces in smart building design~\cite{tschaepe2021somaesthetics, rao2025smart}. Third, we envision \emph{participatory and justice-oriented} directions for HBI. Our opportunities include VR platforms for community co-design and systems that register non-human presence—such as pollinators or wildlife corridors \textbf{(DO3, DO4, DO5)}. These interactions extend critical HCI into architectural contexts~\cite{hirsch2023re}, positioning smart buildings as spaces of care, activism, and more-than-human coexistence.


\subsection{Reflective Messages for Multi-sensory and Multi-dimensional Biophilic Interaction in HBI}
Smart buildings must evolve to support richer, more ecologically attuned forms of interaction, aligned with a multisensory - multidimensional (MS–MD) approach. To support this reorientation, we distil five reflective messages by reflecting on the combined study findings, intended to guide future HBI design and research. First, biophilic interaction should be scaled across different \emph{perceptual levels}---from micro (e.g. tactile encounters through fingertips), to meso (e.g. room-scale soundscapes or kinetic partitions), to macro (e.g. building-wide lighting systems reflecting outdoor conditions). Second, building on this multi-scale sensitivity, interactions should unfold \emph{temporally}, allowing environments to act as narrative spaces shaped by seasonal rhythms and lived memory. Third, biophilic systems should strive for \emph{coherence} across sensory modalities, so that light, sound, temperature, and motion collectively evoke ecological meaning rather than functioning as disconnected features. Fourth, well-timed ``disruptions'' or moments of \emph{surprise} can play a powerful role—fostering curiosity, reflection, and deeper connection to place and self. Finally, biophilic interaction must remain \emph{inclusive}, designed with sensitivity to diverse bodies, access needs, and cultural ways of sensing and being. Together, these messages reframe biophilic interaction as interpretive (inviting meaning-making) rather than prescriptive (dictating responses), and reciprocal rather than extractive---qualities that are critical as HBI moves toward more-than-human, justice-oriented futures.


\subsection{Discussion of Biophilic Interaction Design Tensions}
The opportunities outlined in this paper highlight a critical tension between technological advancement and the intrinsic, cultural, and emotional values of nature~\cite{gunderson2014social}. This ``technocratic fallacy'' reflects the risk of framing biophilic design through an overly technical lens. Such framings may prioritise quantifiable metrics, such as run-time efficiency of IoT systems managing building conditions, side-lining the intrinsic qualities of nature that evoke awe, connection, and belonging. Detached, technical discourse can further disconnect occupants from the authentic and restorative qualities of nature. We need to remain mindful of such framings and ensure technological solutions complement the humanistic and ecological principles of biophilia~\cite{kellert1995biophilia, wilson1986biophilia, kellert2018nature}.

Biophilic design also entails a reciprocal relationship. One side concerns humans benefiting from nature---through sensory stimulation and well-being; the other involves nature benefiting from us---through care, stewardship, and infrastructural support. Expert visions largely focused on the former, envisioning how smart buildings might enhance human experience. This human-centred focus mirrors a broader gap in HBI, where technologies often overlook ecological reciprocity. Yet in the face of climate change, biodiversity loss, and environmental degradation, reciprocity is essential. Smart buildings are not passive---they reshape ecosystems, consume resources, and produce waste. They must therefore be designed not only to evoke nature, but to sustain it. Moving forward, we call for biophilic design in HBI to be critically reoriented toward ecological responsibility and care for more-than-human life.

\subsection{Discussion of Limitations and Future Work}
This work reflects our positional lens as researchers working across HCI, architecture, and interaction design, and is not without limitations. First, most experts were European, and their imaginaries may not reflect broader societal, Indigenous perspectives on biophilic interaction. While we observed disciplinary differences---such as technical versus aesthetic framings---most experts shared a design-forward orientation, potentially skewing the findings toward certain types of interventions. Some offered perspectives in social justice or spirituality, but these remained peripheral. As one noted: \textit{"European architecture focuses on buildings for human design---a lot of this comes from the dominance of Christianity, which emphasised actions in favour of humankind. This contrasts with the perspectives of say the Indigenous tribes in Australia..."}. Future work should intentionally engage more diverse world-views and situated knowledge.

Second, we chose to conduct interviews over scoping or systematic literature reviews. We find that interviews allowed us to capture tacit knowledge, practice-based reflections, and design imaginaries that may not yet be documented in academic discourse. Moreover, the initial coding was conducted by a single author. Although findings were iteratively discussed and refined, this may introduce subjective emphasis. Future work could extend our interview data with methods---such as co-design, literature reviews, or in-situ prototyping---to broaden empirical grounding. Finally, the technological examples linked to our design opportunities were interpretively developed. Although not discussed directly in interviews, we drew on our disciplinary expertise to propose plausible interventions---such as generative AI and ecological sensors---chosen for their conceptual resonance with the experiential qualities described by experts. We offer these as illustrative pathways to expand the HBI design space, and we invite future research to adapt, extend, or critically examine these interventions through prototyping and participatory methods.




 